\chapter{Pupils, diffraction, resolution, and focus}\label{ch:diffraction}
\reading{Saleh--Teich \S\S4.1--4.4, printed pp.~105--137.
Start at PDF pp.~127, 138, 143, and 149.
Born--Wolf \S8.3, printed p.~375 (PDF p.~413), and \S8.5.2,
printed p.~395; the Airy formula is on p.~396 (PDF p.~434).}

\section{Aperture stops, pupils, and numerical aperture}
An aperture stop limits the cone of rays that an optical system accepts
from an object point. The entrance pupil is the image of that stop seen
through the optics before it; the exit pupil is its image through the
optics after it. A pupil can therefore be a virtual image rather than a
physical hole at the plane where it is drawn.

For a rotationally symmetric image-side cone of half-angle $\alpha_{\max}$,
\begin{equation}
\NA=n\sin\alpha_{\max}.
\label{eq:NA}
\end{equation}
In vacuum $n=1$. Thus $\NA=0.33$ corresponds to
$\alpha_{\max}=\arcsin(0.33)\simeq19.27^\circ$.
The chief ray goes through the centre of the aperture stop; marginal rays
reach its edge. These are different rays.
The mask chief-ray angle of $6^\circ$ in Table 1 is not the same angle as
the wafer-side numerical-aperture cone.

In a paraxial system with pupil radius $a$ and effective image distance $f$,
$\NA\simeq a/f$ in air. The exact sine definition is the basic one.
The optical invariant gives approximately
$\NA_{\rm object}\simeq |M|\NA_{\rm image}$ for equal media and lateral
magnification $M$. A 4:1 reduction has $|M|=1/4$; object and image NAs
must not be silently interchanged.

\section{Propagation as a sum of angular components}
Insert a transverse Fourier component
$\exp[\ii2\pi(f_xx+f_yy)]\exp(\ii k_z z)$ into the Helmholtz equation
in vacuum. The result is
\[
k_z^2=k_0^2-(2\pi)^2(f_x^2+f_y^2).
\]
Each propagating component therefore acquires
\begin{equation}
H_z(\vct f)=
\exp\!\left[\ii k_0z\sqrt{1-\lambda^2|\vct f|^2}\right].
\label{eq:angularprop}
\end{equation}
For $|\vct f|>1/\lambda$, the longitudinal wave number is imaginary and
the component is evanescent. The far-field projector does not collect
arbitrarily fine near-field detail.

When $\lambda|\vct f|\ll1$,
\begin{equation}
H_z(\vct f)\simeq
\ee^{\ii k_0z}\ee^{-\ii\pi\lambda z|\vct f|^2}.
\label{eq:fresneltransfer}
\end{equation}
Inverse transformation yields the Fresnel propagation integral
\begin{equation}
U(x,y;z)=\frac{\ee^{\ii k_0z}}{\ii\lambda z}
\iint U(x',y';0)
\exp\!\left[\frac{\ii\pi}{\lambda z}
\{(x-x')^2+(y-y')^2\}\right]\dd x'\dd y'.
\label{eq:fresnel}
\end{equation}
One way to obtain the same phase is to expand the distance
$R=\sqrt{z^2+(x-x')^2+(y-y')^2}$ as
$z+\{(x-x')^2+(y-y')^2\}/(2z)$.
The amplitude $1/R$ is approximated by $1/z$.
These are two approximations, and the phase approximation is often
the stricter one because it is multiplied by the large wave number.

If the quadratic source-coordinate phase is negligible,
$a^2/(\lambda z)\ll1$ for characteristic aperture radius $a$,
Eq.~\eqref{eq:fresnel} becomes a Fourier transform:
\[
U(x,y;z)\propto
\widetilde U\!\left(\frac{x}{\lambda z},\frac{y}{\lambda z}\right)
\]
apart from a known output quadratic phase. This is Fraunhofer diffraction.
The dimensionless ratio $N_F=a^2/(\lambda z)$ is the Fresnel number.
A focusing optic can cancel the required quadratic phase, giving a
Fourier-related focal plane without requiring a physically enormous
propagation distance.

\section{The pupil function and the point-spread function}
Define a scalar coherent transfer function
\begin{equation}
P(\vct f)=A(\vct f)\,
\mathbf1_{|\vct f|\leq f_c}\,
\exp\!\left[\frac{\ii2\pi}{\lambda}W(\vct f/f_c)\right],
\qquad f_c=\frac{\NA}{\lambda}.
\label{eq:pupil}
\end{equation}
Here $\mathbf1$ equals one inside the admitted disk and zero outside,
$A$ describes amplitude transmission or reflection, and $W$ is the
phase-equivalent OPD. The normalized pupil coordinate is
$\vct u=\vct f/f_c=(\rho\cos\theta,\rho\sin\theta)$.

For coherent illumination of a thin scalar object with complex
transmittance or reflectance $t$,
\begin{equation}
U_{\rm image}=\FT^{-1}[\widetilde t\,P]=t*h,\qquad
h=\FT^{-1}[P].
\label{eq:coherentimage}
\end{equation}
$h$ is the amplitude point-spread function (PSF).
The intensity PSF is $|h|^2$. A point object contains all transverse
spatial frequencies; the pupil selects the part used to reconstruct it.
This is why a finite aperture produces a finite spot.

For a clear circular pupil with no aberration, let $r=\sqrt{x^2+y^2}$.
Angular integration gives a Bessel function:
\[
h(r)=2\pi\int_0^{f_c}fJ_0(2\pi fr)\dd f
=\frac{f_cJ_1(2\pi f_cr)}{r}.
\]
Since $h(0)=\pi f_c^2$, the normalized intensity is
\begin{equation}
\frac{I(r)}{I(0)}
=\left[\frac{2J_1(v)}{v}\right]^2,\qquad
v=2\pi f_cr.
\label{eq:airy}
\end{equation}
The first zero of $J_1$ is $v\simeq3.8317$, so
\begin{equation}
r_{\rm first\ zero}=0.6098\,\frac{\lambda}{\NA}.
\end{equation}
This is the Airy pattern. Its rings redistribute energy rather than
indicating geometric rays bouncing around the image.

\section{Coherent and incoherent resolution are different questions}
For coherent illumination, amplitudes from object points interfere before
we square the field. For mutually incoherent object points, cross terms
average out and object \emph{intensity} is convolved with $|h|^2$.
The incoherent optical transfer function is consequently a pupil
autocorrelation:
\begin{equation}
\mathrm{OTF}(\vct q)=
\frac{\int P(\vct f+\vct q/2)P^*(\vct f-\vct q/2)\dd^2f}
{\int |P(\vct f)|^2\dd^2f}.
\label{eq:otf}
\end{equation}
Two disks of radius $f_c$ cease to overlap when their separation exceeds
$2f_c$. This explains the incoherent intensity cutoff $2\NA/\lambda$,
whereas the on-axis coherent field cutoff is $\NA/\lambda$.
Partial coherence, used in lithography, requires the source-dependent
formulation developed later.

The Rayleigh two-point separation $0.61\lambda/\NA$ is a criterion for
distinguishing point-source images. It is not a universal minimum
thresholded hole diameter. Periodic patterns, oblique illumination,
resist thresholds, and proximity correction pose a different problem.
Do not conclude that the paper's 14 nm contacts are impossible merely
because $0.61(13.5)/0.33\simeq24.95\nm$.

\section{Resolution and depth-of-focus scaling}
Lithography often summarizes a process-dependent minimum feature scale as
\begin{equation}
\mathrm{CD}\sim k_1\frac{\lambda}{\NA},
\label{eq:k1}
\end{equation}
where $k_1$ absorbs pattern, illumination, mask, resist, and process choices.
For 14 nm, 13.5 nm wavelength, and 0.33 NA, the numerical ratio is
$k_1=14(0.33)/13.5\simeq0.3422$.
This calculation characterizes the geometry; it does not prove a process
window or assign the same $k_1$ to every pitch.

Defocus adds a relative phase from Eq.~\eqref{eq:angularprop}.
Subtract the on-axis piston:
\begin{align}
\phi_z(\rho)
&=\frac{2\pi z}{\lambda}
\left[\sqrt{1-\NA^2\rho^2}-1\right]\\
&\simeq-\frac{\pi z\NA^2}{\lambda}\rho^2.
\label{eq:defocusphase}
\end{align}
When the pupil-edge phase becomes a fixed tolerable number, the
corresponding defocus scales as $\lambda/\NA^2$:
\begin{equation}
\mathrm{DOF}\sim k_2\frac{\lambda}{\NA^2}.
\end{equation}
Whether DOF denotes a half-range or a full range changes the convention
for $k_2$. The paper specifies sampled focus values, not a universal $k_2$.

At its conditions,
\[
\lambda/\NA=40.91\nm,\qquad
\lambda/\NA^2=123.97\nm.
\]
Increasing NA improves lateral resolution but generally makes focus more
sensitive. Reducing wavelength improves resolution but also reduces this
focus scale at fixed NA. Reducing NA has the opposite resolution effect.
The Introduction's combined discussion of smaller wavelength and NA
should be read with these two distinct scalings in mind.

\section{A quantitative focus check}
From Eq.~\eqref{eq:defocusphase},
\[
W_z(\rho)\simeq-\frac{z\NA^2}{2}\rho^2.
\]
Remove its mean and use $\Var(\rho^2)=1/12$:
\begin{equation}
\sigma_{W,z}=\frac{|z|\NA^2}{4\sqrt3}.
\label{eq:defocusRMS}
\end{equation}
At $z=15\nm$, $\sigma_{W,z}=0.235775\nm
=17.4648\,\mathrm{m}\lambda$.
The pupil-edge-to-centre OPD is $0.81675\nm$.
These two values are different because RMS, peak-to-valley, and
edge-to-centre error are different statistics.

The exact angular formula gives an edge OPD
$15[1-\sqrt{1-0.33^2}]\simeq0.8403\nm$.
The paraxial value is a useful approximation, not an exact
high-angle propagation result.

\chapter{Why EUV needs mirrors and a three-dimensional mask}\label{ch:euv}
\paperloc{Fig.~1, Table 1, Fig.~3, and the mask discussion in Section 2A.}
\reading{Born--Wolf \S1.5, printed p.~36 (PDF p.~65), and
\S1.6, p.~51 (PDF p.~80); its characteristic-matrix discussion begins
at p.~55 (PDF p.~84). Saleh--Teich \S6.2, p.~209 (PDF p.~231).}

\section{The optical path through an EUV exposure system}
The useful sequence is source, collector and illumination optics,
reflective patterned mask, projection mirrors, and resist-coated wafer.
Illumination optics control the angular distribution incident on the
mask. Projection optics determine which mask-diffracted components reach
the wafer and with what amplitude and phase.

At 13.5 nm, strong absorption makes ordinary transmissive lens trains
and propagation through air unsuitable. EUV scanners use vacuum and
multilayer mirrors. ASML's NXE:3400C documentation specifies reflective
4:1 reduction optics and 0.33 NA \cite{asml}.
The paper's Fig.~1 caption says NXT3400; the relevant EUV family is
NXE:3400. This naming issue does not alter the optical calculations.

A geometrical mirror's reflection direction is largely independent of
wavelength, so an all-reflective design avoids ordinary refractive
chromatic aberration. Its coatings still have wavelength-dependent
reflectance and phase. Mirrors also absorb energy, heat, deform, and move.
Thus reflective does not mean perfectly achromatic in every optical
property or immune to thermal aberration.

\section{Complex refractive index and absorption}
With time convention $\ee^{-\ii\omega t}$, write a passive EUV material's
index as $\widetilde n=1-\delta+\ii\beta$, where $\beta>0$.
A normally propagating field behaves as
\[
U(z)=U(0)\ee^{\ii k_0(1-\delta)z}\ee^{-k_0\beta z}.
\]
Hence
\begin{equation}
I(z)=I(0)\ee^{-\mu_{\rm abs}z},\qquad
\mu_{\rm abs}=\frac{4\pi\beta}{\lambda}.
\end{equation}
The real part affects phase; the imaginary part attenuates amplitude.
Material optical constants therefore influence both image contrast and
phase. Treating an absorber as merely a perfectly black two-dimensional
painted region omits this physics.

\section{Reflection at an interface}
Continuity of tangential electric and magnetic fields at an interface
gives, for a chosen polarization, equations of the form
\[
1+r=t,\qquad Y_0(1-r)=Y_1t.
\]
Here $Y_j$ is the appropriate optical admittance for that polarization.
Solving,
\begin{equation}
r_{01}=\frac{Y_0-Y_1}{Y_0+Y_1},\qquad
t_{01}=\frac{2Y_0}{Y_0+Y_1}.
\label{eq:fresnelr}
\end{equation}
For electric field perpendicular to the incidence plane,
$Y_j$ is proportional to $\widetilde n_j\cos\theta_j$.
For the tangential electric-field convention of parallel polarization
it is proportional to $\widetilde n_j/\cos\theta_j$.
Different polarization amplitude conventions can reverse a reflected
amplitude sign; physical intensities and phase comparisons must use
one convention consistently.

Snell's relation is
$\widetilde n_0\sin\theta_0=\widetilde n_1\sin\theta_1$.
For absorbing layers the transmitted angle and longitudinal wave number
are generally complex. The physical branch decays into the medium.

\section{A thin film: summing all internal reflections}
Consider a layer 1 of thickness $d$ between media 0 and 2.
The phase for one transit is
$\delta_1=k_0\widetilde n_1d\cos\theta_1$.
The first reflection is $r_{01}$.
The first wave returning from the lower interface contributes
$t_{01}t_{10}r_{12}\ee^{2\ii\delta_1}$.
Each additional round trip multiplies this by
$r_{10}r_{12}\ee^{2\ii\delta_1}$.
Summing the geometric series,
\[
r_{\rm film}
=r_{01}+
\frac{t_{01}t_{10}r_{12}\ee^{2\ii\delta_1}}
{1-r_{10}r_{12}\ee^{2\ii\delta_1}}.
\]
Using $r_{10}=-r_{01}$ and
$t_{01}t_{10}=1-r_{01}^2$ gives
\begin{equation}
r_{\rm film}=
\frac{r_{01}+r_{12}\ee^{2\ii\delta_1}}
{1+r_{01}r_{12}\ee^{2\ii\delta_1}}.
\label{eq:thinfilm}
\end{equation}
Replace the lower-interface reflection by the effective reflection of
all deeper layers and repeat recursively. This constructs a multilayer
response. The equivalent transfer-matrix method propagates tangential
fields through each layer and multiplies the matrices in physical order.
The reflected intensity is $|r_{\rm stack}|^2$ in the incident medium,
and the reflected phase is $\arg r_{\rm stack}$.

This derivation explains why a small thickness error can change
reflectance and phase simultaneously. It also explains why a multilayer
reflector cannot be represented by surface height alone.

\section{The Mo/Si period and the Bragg estimate}
Reflections from successive layers can reinforce when their round-trip
phase difference is approximately $2\pi$. Neglecting refractive
corrections for an estimate, a repeated period $d$ at angle $\alpha$
from the normal obeys
\begin{equation}
2d\cos\alpha\simeq m\lambda.
\label{eq:bragg}
\end{equation}
This is a design estimate; the full result follows from the complex
multilayer calculation above, including interfaces, absorption, and
termination.

The paper gives Mo thickness $0.725\nm$, Si thickness $1\nm$, and
absorber height $15\nm$, explicitly in \emph{wafer scale}. Under a 4:1
scale conversion these correspond to physical mask values
$2.9\nm$, $4\nm$, and $60\nm$.
The inferred period is $6.9\nm$, and
\[
2(6.9)\cos6^\circ=13.7244\nm,
\]
which has the expected scale for 13.5 nm multilayer reflection.
This is a consistency check based on a 4:1 interpretation, not a
recovered complete mask prescription.

\careful{Do not put the wafer-scaled 1.725 nm period directly into a
physical 13.5 nm mask-stack solver. Lateral pattern dimensions can be
reported at wafer scale for convenience, but a Maxwell calculation must
use a consistent physical mask geometry, wavelength, and reduction mapping.
The paper does not provide all stack details needed for exact reconstruction.}

\section{Finite absorber height breaks the ideal thin-mask picture}
An obliquely incident ray traversing absorber height $h$ has a lateral
geometrical displacement $h\tan\alpha$. With the inferred
$h=60\nm$ and $\alpha=6^\circ$, this is about $6.306\nm$ on the mask,
or $1.577\nm$ after 4:1 reduction.
This shows why an apparently small angle matters on nanometre scales.
It is not a prediction of the paper's PS: coherent diffraction, material
phase, double-pass geometry, and contour formation also contribute.

Different diffraction orders interact with the finite absorber and
multilayer differently. Their complex amplitudes can depend on source
angle, polarization, hole orientation, and pitch. This produces
asymmetric images, orientation-dependent CDs, and focus-dependent
placement. A square mask opening can therefore have unequal printed
horizontal and vertical dimensions.

The mask chief-ray angle describes the central illumination geometry.
The quasar distribution describes angular offsets around that geometry.
Off-axis source shaping and a nonzero chief-ray angle are distinct ideas;
one should not infer that the chief ray itself supplies every resolution
benefit attributed to source shaping.

\section{What a rigorous mask calculation has to do}
A representative electromagnetic workflow is:
\begin{enumerate}
\item Specify each material's complex optical constants, layer thickness,
feature shape, sidewall angle, period, source angle, and polarization.
\item Solve the frequency-domain Maxwell boundary-value problem. Methods
include Fourier modal techniques and finite-element discretizations.
\item Obtain the complex vector amplitudes of reflected diffraction orders.
\item Propagate those orders through the projection system, including pupil
phase, amplitude, and defocus.
\item Sum coherent orders for one source point, average intensities over
mutually incoherent source points, and extract image or resist contours.
\end{enumerate}
For a periodic permittivity
$\epsilon(x,z)=\sum_g\epsilon_g(z)\ee^{\ii gx}$, expanding the field in
the same harmonics turns multiplication by $\epsilon$ into convolution
of Fourier coefficients. The diffraction orders become coupled equations,
not independent rays. Their boundary conditions enforce the material
interfaces. This is the basic reason rigorous diffraction is more
computationally expensive than a binary thin-mask Fourier transform.

The paper identifies a commercial Synopsys simulator but does not give
enough settings to determine its exact electromagnetic solver, numerical
convergence, material database, or resist extraction procedure.
Those items remain unknown rather than being supplied by the word rigorous.
